\subsection{Introduction}
\begin{Note}
Look into VABS literature!
\end{Note}
% The purpose of this section is to establish a systematic framework for deriving 
% six-dimensional section constitutive laws from the Saint-Venant solution class. 
In this section we identify a minimal set of cross sectional warping shapes \(\bm{\Phi}(\bm{r})\) such that the representation \eqref{eq:trace-lift-linear}:
\RepeatEquation{eq:trace-lift-linear}
% \begin{equation*}
% % \begin{SaveEquation}{eq:trace-lift-linear}
% \hat{\bm u}(\reflenx, \bm{r})=\bm{u}(x)
% +\bm\theta(x)\times\bm{r}+\bm\Phi(\bm{r})\,\bm{\alpha}(x)
% % \label{eq:trace-lift-linear}
% % \end{SaveEquation}
% \end{equation*}
holds identically for any solution \(\hat{\bm u}\) to a problem in the Saint-Venant class. 
The representation \eqref{eq:trace-lift-linear} is non-unique, which
poses a fundamental obstacle for associating these modes to classical \emph{deflection} \(\bm{u}(\reflenx)\), \emph{rotation} \(\bm{\theta}(\reflenx)\), and \emph{warping} \(\bm{\alpha}(\reflenx)\) fields. 
The practical implications of this redundancy are substantial, and manifest through several important characteristics:
\begin{description}
    \item[Boundary Conditions] The Saint Venant class prohibits true Dirichlet boundaries. This gives rise to spring-like behavior at fixed ends which is easy to misinterpret. 
    % This is phenomenon is elaborated in \cref{rem:boundary-springs}.
    \item[Coupling] 
    \item[Symmetry] If the strains \(\bm\gamma,\bm\kappa\) generated by the trace are not work-conjugate to the classical resultants \(\bm{n},\bm{m}\) then the model could give rise to an asymmetric tangent stiffness in a finite element simulation. 
    \item[Compatibility] Can the trace be represented by a classical 6-DOF element.
\end{description}
Thus, different identifications of \(\bm{u}\) and \(\bm{\theta}\) may generate distinct  strain measures 
$(\bm{\gamma}, \bm{\kappa})$, constitutive couplings, and boundary conditions, but yet still represent the same particular Saint-Venant solution exactly. 
The ambiguous nature of Saint Venants solutions is well known in the continuum mechanics literature. 
However, we are not aware of any works that formalize this in the context of dimensionally reduced 1D beam theories, or systematically investigate the extent to which beam theories may exactly represent 3D solutions in the Saint Venant class. 
Rather, it is typically assumed at the outset that the 1D model will only approximate a continuum solution. 
The subject is raised by \cite{simo1982bifurcation}, where a brief investigation is made in the context of planar buckling. 
However, because this investigation was limited to the plane, several salient three dimensional considerations do not arise such as work-conjugacy of resultants and handling of asymmetric sections. 



The trace abstraction represents several prominent shear-deformable beam theories:
\begin{description}
    \item[Classical] \cite{timoshenko1987theory} investigated a the flexural solution identified by the root boundary condition:
    \[
    \hat{\bm u}'(x_0, \mathbf{0}) = \mathbf{0}
    \]
    A trace defined to realize boundary conditions in this form exactly reproduces the response of classical Euler-Bernoulli beam theory, while remaining exact in the Saint-Venant class.

    \item[Tangent] \cite{timoshenko1921correction} imposed the root condition by requiring the directional derivative of the axial displacement component along a chosen cross-section direction \(\mathbf{t}\) vanishes at that boundary point. 
    In 3D this might be stated as:
    $$
    \left.\mathbf{t} \cdot \nabla_r(\hat{\bm u} \cdot \FrameBasis)\right|_{x=L, \bm{r}=r_0}=0,
    $$
    This produces a shear correction of \(\kappa = 2/3\).
    % The full trace is derived in \cref{sec:trace-timo}.
    \item[Geometric]  In the landmark work by \cite{cowper1966shear}, a systematic investigation was conducted to obtain shear-correction factors corresponding to a trace which we have generalized to 3D with the form:
    \begin{equation}
    \bar{\bm{u}}(x) \coloneq \SectionAverage{\hat{\bm{u}}} - \bar{\bm{\theta}} \times \CentroidLocation, 
    \qquad 
    \bar{\bm{\theta}}(x) \coloneq \RotationAverage{\hat{\bm{u}}},
    \label{eq:def-section-avg-trace}
    \end{equation}
    Here $\SectionAverage{\cdot}$ and $\RotationAverage{\cdot}$ are integral functionals defined by \eqref{eq:def-section-averages} which average the 3D field \(\hat{\bm u}\) at a cross section and furnish an $L^2$-optimal fitting of a section-rigid field to the 3D displacement \(\hat{\bm u}_{\rm P}\). 
    The free parameters are determined by section geometry through rotation and translation averages of the warping functions, which is carried out in \cref{sec:trace-cowper}.
    The approach by \cite{cowper1966shear} laid the foundation for the plane beam theories developed by \cite{simo1982bifurcation,hjelmstad1983seismic,hjelmstad1987warping}. 

    \item[Conjugate] The shear-correction factors of \cite{krenk1985pretwist}, \cite{renton1991generalized} and \cite{romano1992shear} are realized by a trace that defines beam strains to be 
    work-conjugate to the classical resultants:
    \begin{equation}
    \delta\bm{\gamma} \cdot \bm{n} + \delta\bm{\kappa} \cdot \bm{m} 
    = \int_\Section \delta\bm{\varepsilon} \cdot \bm{\sigma} \, \mathrm{d}A
    \label{eq:def-trace-conjugate}
    \end{equation}
    for all variations within the Saint~Venant class.
      % \begin{equation*}
      %   \int \ShearStress_{(b)}\cdot (\ShearStress_{(b)}/G)\,\mathrm{d}A
      %   = \ShearForce\cdot \bm\gamma^{\rm sh} + T_{\mathrm C^{\rm sh}}\varkappa_{(b)}
      % \end{equation*}
    This uniquely determines all 18 
    parameters and guarantees a symmetric section stiffness matrix. The full trace is derived in \cref{sec:trace-energy}. 
    \item[Alternating]
    \cite[Eq. (7.4.52)]{slaughter2013linearized}, and \cite{reddy2002energy}:
    \[
    \tfrac12 \bm{\epsilon}\nabla\times\hat{\bm u}= \mathbf{0}
    \]
    This produces a \emph{stiffening} shear correction of \(\kappa = 4/3\), which is peculiarly greater than unity. 
    % The full trace is derived in \cref{sec:trace-slaughter}.
\end{description}

\subsubsection{Literature}
% In the literature, warping shapes have been identified by various approaches:
The most appropriate warping shape for torsion is the Saint~Venant function defined by \eqref{eq:bvp-warp-twist-iesan}. 
This works well in Type 1, 2, and 3 formulations as several special cases of Saint~Venant's torsion problem can be exactly represented in a Type 3 element by associating this shape directly to the twist \(\Twist\).
The case of flexural (shear) warping is considerably more complicated, and the Saint~Venant flexural warpage \(\WarpShearIesan\) defined by \eqref{eq:bvp-warp-shear-iesan} does not clearly associate with a classical deformation measure \(\bm{\gamma}\) or \(\bm{\kappa}\). 
In a Type 1 formulation where auxiliary degrees of freedom are permitted:
\begin{itemize}
    \item \cite{lecorvec2012nonlinear} employed a general interpolation basis. This is a very general approach, but generally requires several modes to achieve reasonable accuracy.
    \item  \cite{klinkel2003anisotropic} employed the classical Saint~Venant warping functions defined by \eqref{eq:bvp-warp-shear-iesan}. 
\end{itemize}


\begin{Prompt}
We currently have two separate lines of thought: the beam kinematic postulate $\bm{u} + \bm{\theta}\times \bm{r} + \bm{\Phi} \bm{\alpha}$, which is very general, and the trace strain matrix $T_0 \bm{p}$ which pulls us inescapably into Saint Venant land. 
How could we have arrived at this same point directly from the beam kinematic decomposition? what lies at the heart of what we've done? 
By saying $\bm{u} + \bm{\theta}\times \bm{r} + \bm\Phi \bm\alpha$ is ambiguous, really what were saying is that boundary conditions on $\bm{u},\bm{\theta}$ do not have to necessarily be boundary conditions on the plane-rigid basis, if we absorb part of this basis into $\bm{\Phi}$. 
This is the practical perspective, because specified boundary conditions in the reduced trace fields are the mechanism through which the framework produces a beam theory with different behavior. 
So say that its now $\bm{u}(x) + \bm{\theta}(x)\times \bm{r} + \bm{\Phi} \bm{\alpha}(x) + \bm{\Psi} \bm{\eta}(x)$, where $\bm{\Psi}$ is the plane-rigid basis. 
Recall that one of the things we're after is a general expression for the material strains that is not so dependent on the load parameters, at least not in the plane-section basis... they are inescapable in the pure-warping basis because those warping functions naturally feature them by construction.
\end{Prompt}

\begin{Prompt}
My objective is to develop a finite element formulation which is capable of representing a wide variety of traces. my hypothesis is that many of these will be useful in practice, and could provide modelers with an additional kinematic freedom to control  boundary conditions. Typically this subject is approached with the objective of finding *one* ideal trace, which i believe could be excessively restrictive. 
I believe such an approach is overly dogmatic, and fails to recognize the utility that other traces offer. 
I've never seen anyone discuss the subject with my general outlook; I've also never seen anyone acknowledge the possibility there may exist multiple valid traces. 

I am aware of a plethora of literature that attacks the question: ``under what conditions are lower dimensional beam theories valid approximations to continuum mechanics". These works tend to fixate on *one* classical beam theory (ie, classical euler or timoshenko) and ask under what conditions they're valid. 
What I'm doing is different; my work is answering: ``if we loosely generalize the idea of a beam theory, while maintaining constraints which ensure practicality, to what extent can we \emph{exactly} represent a given class of continuum solutions".
\end{Prompt}
